\section{系统的复频域分析}\label{sec:Complex_Freq_Analysis}

\subsection{零极点图与频率响应特性}\label{sub:zero_pole_characteristics}

对于最典型的因果线性时不变系统，其系统函数为有理函数，可以将系统函数表示为：
\begin{align*}
    H(s)=K\frac{(s-z_1)(s-z_2)\cdots (s-z_m)}{(s-p_1)(s-p_2)\cdots (s-p_n)}
\end{align*}
其中$z_1,z_2,\cdots ,z_m$为系统函数的零点，$p_1,p_2,\cdots ,p_n$为系统函数的极点，$m<n$，$K$为常数。将零点和极点在复平面上表示出来，得到\textbf{零极点图}(zero-pole plot)，其中，零点用“O”表示，极点用“X”表示，高阶的零点或极点，在其旁边用数字标明其重数。下面给出一个例子：

\begin{example}
    设系统函数为
    \[H(s)=\frac{s+2}{(s+1)^2 (s+3)}\]
    则其零极点图如下所示：
    \begin{figure}[H]
        \centering
        \includegraphics[width=0.4\textwidth]{zero_pole}
        \caption{零极点图(1)}
    \end{figure}
\end{example}

根据\ref{sub:stability_of_system}中提出的系统的稳定性条件，因果的线性时不变系统是稳定的，当且仅当其系统函数的收敛域包含虚轴，此时系统函数的所有极点均位于左半平面$\text{Re}\ s<0$，\textbf{系统的频率响应特性即系统函数在虚轴上的限制}。

此时，如果零极点图较为简单，则可以由零极点图判断系统的滤波特性，即判断系统函数在虚轴上不同位置的幅值大小。

零点通常会使得附近的系统函数幅值减小，特别地，如果虚轴上存在零点，则可直接判断该点处的幅频特性值为0：

\begin{example}
    设系统的零极点图如下图(a)所示：
    \begin{figure}[H]
        \centering
        \includegraphics[width=0.7\textwidth]{zero_pole1}
        \caption{零极点图与幅频特性曲线(1)}
    \end{figure}
    可以看到系统函数在原点处由一个一阶零点，这表明系统的频率响应特性在$\omega=0$有零点，也就意味着系统对低频信号有抑制作用；系统的极点均位于实轴上，这意味着系统函数的分母是有两个实根的实系数多项式
    \[H(s)=\frac{Ks}{(s+a)(s+b)}\]
    考虑它在虚轴上的限制的幅值：
    \[|H(\mathrm{i}\omega)|=\frac{K\omega}{\sqrt{(\omega^2+a^2)(\omega^2+b^2)}}\]
    当$\omega\to\infty$时，$|H(\mathrm{i}\omega)|\to 0$，这表明系统对高频信号有抑制作用。因此，该系统具有带通滤波特性。

    可以定性地画出其幅频特性曲线，如上图(b)所示。
\end{example}

极点通常会使得附近的系统函数幅值增大：
\begin{example}
    设系统的零极点图如图下图(a)所示：
    \begin{figure}[H]
        \centering
        \includegraphics[width=0.7\textwidth]{zero_pole2}
        \caption{零极点图与幅频特性曲线(2)}
    \end{figure}

    系统函数的分母在虚轴左侧有一对共轭复根，这意味着系统函数形如
    \[H(s)=\frac{K}{(s-p)(s-\bar{p})}=\frac{K}{s^2-2\text{Re}(p)s+|p|^2}\]
    考虑它在虚轴上的限制的幅值：
    \[|H(\mathrm{i}\omega)|=\frac{K}{\sqrt{(\omega^2-2\text{Re}(p)\omega+|p|^2)}}\]
    由于$\text{Re}(p)<0$，因此容易求得幅频特性在$\omega=\pm \text{Im}(p)$处取得极大值，这表明系统对该频率附近的信号有增强作用。因此，该系统具有带通滤波特性。
    
    可以定性地画出其幅频特性曲线，如上图(b)所示。
\end{example}

以上推断系统函数的形式的过程，可以进一步简化。对于给定位置的零点$z$，其系统函数分子必然包含$s-z$这一因子，从而幅频特性是系统函数的其他部分与一个\textbf{起点为$z$，终点为$\mathrm{i}\omega$的向量长度}的积；对于给定位置的极点$p$，其系统函数分分母必然包含$s-p$这一因子，从而幅频特性是系统函数的其他部分与一个\textbf{起点为$p$，终点为$\mathrm{i}\omega$的向量长度}的商。

\begin{example}
    设系统的零极点图如下图(a)所示：
    \begin{figure}[H]
        \centering
        \includegraphics[width=0.7\textwidth]{zero_pole_vector}
        \caption{零极点图与幅频特性曲线(3)}
    \end{figure}

    考虑不同角频率处的各向量长度，当角频率较小时，各向量的长度都较小；随角频率增大，极点向量的长度增长较快，且极点有两个，因此系统的幅频特性大小随角频率增大而减小，系统具有低通滤波特性。
    
    可以定性地画出其幅频特性曲线，如上图(b)所示。虽然不太容易看出幅频特性曲线在低频段内有一个递增的区间，但这并不影响系统的低通滤波特性。
\end{example}

\subsection{系统流图与梅森增益公式}\label{sub:mason_gain_formula}

本小节讨论多个系统进行组合时的表现，信号流图是研究此问题的有力工具。系统的\textbf{信号流图}(signal flow graph)由香农在1942年发明，后经梅森(J.E. Mason)推广完善，成为分析线性时不变系统的有力工具，因此又称为\textbf{梅森信号流图}(Mason signal flow graph)或\textbf{梅森图}(Mason graph)。

信号流图可以看作是系统框图的另一种表示形式。由于系统框图不能很好地展示出系统的数学模型，本书中未涉及系统框图的概念，感兴趣的读者可以参考\cite{oppenheim1997signals}\cite{signal_systems}。

\begin{definition}
信号流图是一种有向图，其节点表示系统的变量（输入、输出或中间变量），有向边表示变量之间的线性关系，边上的权值表示线性关系的系数。信号流图可以直观地反映系统的结构特性，并且可以通过图论的方法来分析系统的传递函数。

信号流图相关术语汇总如下：\begin{itemize}
    \item \textbf{节点}：表示系统的输入/输出或中间变量，分为\textbf{输入节点（源点）}、\textbf{输出节点（阱点）}和\textbf{混合节点}；
    \item \textbf{支路}：连接两个节点的有向边，表示系统，边上的权值为系统函数，又称\textbf{增益}；
    \item \textbf{通路}：沿支路方向连接的节点序列，其中，不自交的通路称为\textbf{开通路}，从源点到阱点的开通路称为\textbf{前向通路}。将通路上各支路增益之积称为\textbf{通路增益}；
    \item \textbf{环路/闭通路}：起点和终点相同且不自交的通路。将环路上各支路增益之积称为\textbf{环路增益}；
    \item \textbf{不接触环路}：没有公共节点的两环路称为不接触环路；
\end{itemize}
\end{definition}

信号流图与电路有一定的相似之处。两个系统“串联”时，实际上是将两个系统级联，完整系统的系统函数为两个子系统的系统函数之积，乘法交换律告诉我们\textbf{线性时不变系统总是可交换的}；两个系统“并联”时，完整系统的系统函数为两个子系统的系统函数之和（因为激励独立地经过两个系统再在输出端相加）。

\begin{figure}[H]
    \centering
    \begin{minipage}{0.45\textwidth}
        \centering
        \begin{tikzpicture}[auto, >=Stealth, node distance=1.5cm]
            \tikzset{signode/.style={circle, draw, minimum size=2mm, inner sep=0pt, line width=0.6pt}}
            \node[signode, label=left:$x(t)$] (x) {};
            \node[signode, right=of x, label=below:$w(t)$] (w) {};
            \node[signode, right=of w, label=right:$y(t)$] (y) {};
            \draw[->] (x) -- node {$H_1(s)$} (w);
            \draw[->] (w) -- node {$H_2(s)$} (y);
        \end{tikzpicture}
        \caption{串联}
    \end{minipage}
    \hfill
    \begin{minipage}{0.45\textwidth}
        \centering
        \begin{tikzpicture}[auto, >=Stealth, node distance=2.5cm]
            \tikzset{signode/.style={circle, draw, minimum size=2mm, inner sep=0pt, line width=0.6pt}}
            \node[signode, label=left:$x(t)$] (x) {};
            \node[signode, right=of x, label=right:$y(t)$] (y) {};
            \draw[->] (x) to[bend left=45] node {$H_1(s)$} (y);
            \draw[->] (x) to[bend right=45] node[swap] {$H_2(s)$} (y);
        \end{tikzpicture}
        \caption{并联}
    \end{minipage}
\end{figure}

此外，还可以引入反馈回路，根据是否有反馈回路，系统分为开环系统和闭环系统。
\begin{definition}[开环系统与闭环系统]
    如果系统取后级响应为反馈信号输入到前级，则该系统称为\textbf{闭环系统}，否则称为\textbf{开环系统}。
\end{definition}

\begin{figure}[H]
    \centering
    \begin{tikzpicture}[auto, >=Stealth, node distance=1.5cm]
        \tikzset{signode/.style={circle, draw, minimum size=2mm, inner sep=0pt, line width=0.6pt}}
        
        \node[signode, label=left:$x(t)$] (x) {};
        \node[signode, right=of x] (e) {};
        \node[signode, right=of e] (pre_y) {};
        \node[signode, right=of pre_y, label=right:$y(t)$] (y) {};
        
        \draw[->] (x) -- (e);
        \draw[->] (e) -- node {$H_1(s)$} (pre_y);
        \draw[->] (pre_y) -- (y);
        \draw[->] (pre_y) to[bend left=50] node {$H_2(s)$} (e);
    \end{tikzpicture}
    \caption{闭环系统}
\end{figure}

上图展示了一个最简单的闭环系统的信号流图。我们将$H_1(s)$称为\textbf{开环增益}，表示不考虑反馈回路时系统的增益，与之相对，将考虑反馈回路后的增益称为\textbf{闭环增益}。

计算该系统的闭环增益：
\begin{align*}
    Y(s)&=H_1(s)E(s)\\
    E(s)&=X(s)-H_2(s)Y(s)\\
    \implies Y(s)&=H_1(s)\left[X(s)-H_2(s)Y(s)\right]\\
    \implies Y(s)&=\frac{H_1(s)}{1+H_1(s)H_2(s)}X(s)\\
    \implies H(s)&=\frac{Y(s)}{X(s)}=\frac{H_1(s)}{1+H_1(s)H_2(s)}
\end{align*}

\begin{definition}[反馈深度与反馈类型]
将$1+H_1(s)H_2(s)$称为\textbf{反馈深度}，它是引入反馈回路导致的增益下降的倍率。

如果反馈深度大于1，则系统的输出更加稳定，但增益下降，此时反馈为\textbf{负反馈}，特别地，如果反馈深度远大于1，称为\textbf{深度负反馈}，此时$1+H_1(s)H_2(s)\approx H_1(s)H_2(s)$，增益近似等于$1/H_2(s)$（即反馈回路系统的逆系统），与开环增益无关。

如果反馈深度小于1，则系统的输出不稳定，但增益变高，称为\textbf{正反馈}，多用于振荡器；如果反馈深度近似为0，系统将进入\textbf{自激振荡状态}，通常电路无法正常工作，需要消除这种状态。
\end{definition}

\begin{example}
    设$X_1,X_2$为输入节点，$Y_1,Y_2$为输出节点，其他节点为混合节点。
    \begin{figure}[H]
        \centering
        \begin{tikzpicture}[auto, >=Stealth, node distance=1.8cm, every node/.style={font=\small}]
            % node styles: hollow smaller circles
            \tikzset{
                signode/.style={circle, draw, minimum size=2mm, inner sep=0pt, line width=0.6pt},
                lab/.style={font=\small}
            }
            % nodes (empty hollow circles)
            \node[signode] (x1) {};
            \node[signode, below=1.2cm of x1] (x2) {};
            \node[signode, right=2.2cm of x1] (m1) {};
            \node[signode, right=2.2cm of x2] (m2) {};
            \node[signode, right=2.2cm of m1] (y1) {};
            \node[signode, right=2.2cm of m2] (y2) {};
            % external labels (placed to avoid overlap)
            \node[lab, left=1mm of x1] {$X_1$};
            \node[lab, left=1mm of x2] {$X_2$};
            \node[lab, right=1mm of y1] {$Y_1$};
            \node[lab, right=1mm of y2] {$Y_2$};
            % forward path gains use H_i, feedback/return gains use G_j
            \draw[->] (x1) to node[midway, above,  inner sep=1pt] {$H_1$} (m1);
            \draw[->] (x2) to node[midway, below,  inner sep=1pt] {$H_2$} (m1);
            \draw[->] (m1) to node[midway, above,  inner sep=1pt] {$H_3$} (y1);
            \draw[->] (m2) to node[midway, below,  inner sep=1pt] {$H_4$} (y2);
            % adjusted label positions for G_5 and G_6 to avoid overlap
            \draw[->] (m1) to[bend right=8] node[midway, right,xshift=6pt, yshift=4pt,  inner sep=1pt] {$G_5$} (m2);
            \draw[->] (m2) to[bend right=22] node[midway, left, xshift=-6pt,yshift=-4pt,  inner sep=1pt] {$G_6$} (m1);
            \draw[->] (m1) edge[loop above, looseness=50] node[  inner sep=1pt] {$L_1$} (m1);
        \end{tikzpicture}
    \end{figure}
    记上、下两个混合节点对应的中间变量分别为$M_1,M_2$，则
    \begin{align*}
        M_1&=H_1 X_1+H_2 X_2+L_1 M_1+G_5 M_2\\
        M_2&=G_6 M_1\\
        Y_1&=H_3 M_1\\
        Y_2&=H_4 M_2
    \end{align*}
    这是一个以$M_1,M_2,Y_1,Y_2$为未知数的线性方程组，可以解得\begin{align*}
        M_1&=\frac{H_1 X_1+H_2 X_2}{1-L_1-G_5 G_6}\\
        M_2&=\frac{G_6 (H_1 X_1+H_2 X_2)}{1-L_1-G_5 G_6}\\
        Y_1&=\frac{H_3 (H_1 X_1+H_2 X_2)}{1-L_1-G_5 G_6}\\
        Y_2&=\frac{H_4 G_6 (H_1 X_1+H_2 X_2)}{1-L_1-G_5 G_6}
    \end{align*}
\end{example}

信号流图是一种有向图，$n$个节点的信号流图可以通过$n$阶的\textbf{邻接矩阵}$A$来表示，其元素$a_{ij}$表示从节点$j$到节点$\mathrm{i}$的支路增益，如果没有支路连接这两个节点，则$a_{ij}=0$。%显然，信号流图可由其邻接矩阵唯一表示，但邻接矩阵中会有一些冗余的行和列。我们通过一个例子来说明这一点。

\begin{example}
    上例中，节点依次为$X_1,X_2,M_1,M_2,Y_1,Y_2$，则对应的邻接矩阵为
    \[A= \begin{bmatrix}
        0 & 0 & 0 & 0 & 0 & 0 \\
        0 & 0 & 0 & 0 & 0 & 0 \\
        H_1 & H_2 & L_1 & G_5 & 0 & 0 \\
        0 & 0 & G_6 & 0 & 0 & 0 \\
        0 & 0 & H_3 & 0 & 0 & 0 \\
        0 & 0 & 0 & H_4 & 0 & 0
    \end{bmatrix}\]
    在信号流图的语境下，源点对应的行和阱点对应的列均为0。有冗余的行和列，不是我们希望看到的。

    事实上，我们仅仅需要中间节点对应的邻接矩阵。由于响应$Y_1,Y_2$直接取自中间节点，只需要求解$M_1,M_2$，就可获取系统的全部信息。将信号流图表示为一个线性方程组：
    \[\begin{bmatrix}
        M_1 \\
        M_2
    \end{bmatrix}=
    \begin{bmatrix}
        L_1 & G_5\\
        G_6 & 0
    \end{bmatrix}
    \begin{bmatrix}
        M_1 \\
        M_2
    \end{bmatrix}+\begin{bmatrix}
        H_1 X_1 + H_2 X_2 \\
        0
    \end{bmatrix},\text{即}\mathbf{M}=A\mathbf{M}+\mathbf{X}\]
    其中$$A=\begin{bmatrix}
        L_1 & G_5\\
        G_6 & 0
    \end{bmatrix}\text{为中间节点的邻接矩阵},\mathbf{M}=\begin{bmatrix}
        M_1 \\
        M_2
    \end{bmatrix},\mathbf{X}=\begin{bmatrix}
        H_1 X_1 + H_2 X_2 \\
        0
    \end{bmatrix}$$
    整理得\[\begin{bmatrix}
        H_1 X_1 + H_2 X_2 \\
        0
    \end{bmatrix}=\begin{bmatrix}
        1 -L_1 & -G_5\\
        -G_6 & 1
    \end{bmatrix}\begin{bmatrix}
        M_1 \\
        M_2
    \end{bmatrix},\text{即}\mathbf{X}=(I-A)\mathbf{M}\]
    只要矩阵\[S=I-A=\begin{bmatrix}
        1 -L_1 & -G_5\\
        -G_6 & 1
    \end{bmatrix}\]
    可逆，就可以求解出$M_1,M_2$，我们将这个矩阵称为\textbf{系统矩阵}(system matrix)，而将$A$称为\textbf{增益矩阵}。系统矩阵的行列式称为\textbf{特征行列式}，记为$\Delta$。在本例中，有\begin{align*}
        \Delta&=\begin{vmatrix}
        1 -L_1 & -G_5 \\
        -G_6 & 1
    \end{vmatrix}=1-L_1-G_5 G_6
    \end{align*}
    要使系统稳定，必须有$\Delta=1-L_1-G_5 G_6\neq 0$，否则$M_1,M_2,Y_1,Y_2\to\infty$。
\end{example}

现在来考虑更一般的情况：对于有n个中间节点$M_1,M_2,\cdots,M_n$的系统，不妨设激励接入的节点为$M_1$（否则可用叠加原理单独分析各个激励），则系统对应的线性方程组为\begin{align*}
    M_1&=a_{11} M_1 + a_{12} M_2 + \cdots + a_{1n} M_n + X\\
    M_2&=a_{21} M_1 + a_{22} M_2 + \cdots + a_{2n} M_n\\
    &\vdots\\
    M_n&=a_{n1} M_1 + a_{n2} M_2 + \cdots + a_{nn} M_n
\end{align*}
写成矩阵形式为
\[\mathbf{M}=A\mathbf{M}+\mathbf{X}\]
其中$$\mathbf{M}=\begin{bmatrix}
    M_1 \\
    M_2 \\
    \vdots \\
    M_n
\end{bmatrix},\mathbf{X}=\begin{bmatrix}
    X \\
    0 \\
    \vdots \\
    0
\end{bmatrix}$$
在系统矩阵$S=I-A$可逆时，有
\[\mathbf{M}=(I-A)^{-1}\mathbf{X}=S^{-1}\mathbf{X}\]
因此中间节点$M_i$可由以下公式计算：\[M_i=\frac{S^*_{1i}}{\Delta} X\]
其中$\Delta=\det S$为特征行列式，$S^*_{1i}$为系统矩阵$S$的去掉第1行、第i列所得的代数余子式。输出节点的响应是中间节点线性组合。以$M_n$为输出为例（否则可更换节点命名并重复此过程），可以写出\textbf{梅森增益公式}(Mason gain formula)：
\begin{theorem}[梅森增益公式]
    设系统有$n$个中间节点$M_1,M_2,\cdots,M_n$，激励$X$接入节点$M_1$，响应$Y$取自节点$M_n$，则
    \[H=\frac{Y}{X}=(S^{-1})_{1n}=\frac{S^*_{n1}}{\Delta}\]
    其中$(S^{-1})_{1n}$为系统矩阵$S$的逆矩阵的第1行第$n$列元素，$\Delta=\det S$为特征行列式，$S^*_{n1}$为系统矩阵$S$的去掉第$n$行、第1列所得的代数余子式。
\end{theorem}

根据行列式的图论展开的几何意义，梅森增益公式又可以写成如下形式：
\begin{theorem}[梅森增益公式，图论形式]
    设系统有$n$个中间节点$M_1,M_2,\cdots,M_n$，激励$X$接入节点$M_1$，响应$Y$取自节点$M_n$，则
    \[H=\frac{Y}{X}=\frac{\sum_{k}g_k\Delta_k}{\Delta}\]
    其中$g_k$为从源点到阱点的第$k$条前向通路增益，$\Delta_k$为去掉与第$k$条前向通路接触的所有节点后的特征行列式。

    特征行列式可由以下公式计算：
    $$\Delta=1-\sum_a L_a+\sum_{b,c}L_b L_c-\sum_{d,e,f}L_d L_e L_f+\cdots$$
    其中$\sum_a L_a$表示所有环路增益之和，$\sum_{b,c}L_b L_c$表示所有两两不接触环路增益的乘积之和，$\sum_{d,e,f}L_d L_e L_f$表示所有三者互不接触环路增益的乘积之和，依此类推。
\end{theorem}

下面来解释为什么行列式展开具有几何意义。首先，根据行列式的公式：
\[\det S=\sum_{\sigma\in S_n}\text{sgn}(\sigma)\prod_{i=1}^{n}s_{i,\sigma(i)}\]
其中$S_n$为所有$n$个元素的排列组成的集合，$\text{sgn}(\sigma)$为排列$\sigma$的符号。对于$S=I-A$，对角元为$s_{ii}=1-a_{ii}$，非对角元为$s_{ij}=-a_{ij}$。

我们知道，每个置换都可以分解为不相交循环的乘积，例如，$$\sigma(1)=3,\sigma(3)=4,\sigma(4)=1,\sigma(2)=5,\sigma(5)=2$$
可以表示为
$$\sigma=(1\ 3\ 4)(2\ 5)$$
对应中间节点组成的环路
$$M_1\to M_3\to M_4\to M_1,M_2\to M_5\to M_2$$
因此，行列式的每一项$\prod_{i=1}^{n}s_{i,\sigma(i)}$都可以看作是由若干个不接触环路增益的乘积。如果环路中出现不动点$\sigma(\mathrm{i})=\mathrm{i}$，则对应的$s_{ii}=1-a_{ii}$中的1，表示该节点不参与任何环路；如果考虑自环，则对应$-a_{ii}$。

其次，考虑正负号的问题。对于k-循环，它对$\text{sgn}(\sigma)$的贡献为$(-1)^{k-1}$，而每个增益对应一个$-a_{ij}$，因此k-循环对行列式的贡献为$(-1)^{k-1}(-1)^k=-1$。因此，$l$个互不接触的环路增益之积的符号应为$(-1)^l$。
\begin{example}
    考虑具有六个中间节点的信号流图，我们来计算$\sigma=(1)(2)(3\ 4\ 5\ 6)$对应的行列式项：
    \[\text{sgn}(\sigma)\prod_{i=1}^{6}s_{i,\sigma(i)}=s_{1,1}s_{2,2}s_{3,4}s_{4,5}s_{5,6}s_{6,3}=(1-a_{11})(1-a_{22})a_{34}a_{45}a_{56}a_{63}\]
    将因式$(1-a_{11})(1-a_{22})$展开，得到
    \[a_{34}a_{45}a_{56}a_{63}-a_{11}\cdot a_{34}a_{45}a_{56}a_{63}-a_{22}\cdot a_{34}a_{45}a_{56}a_{63}+a_{11}\cdot a_{22}\cdot a_{34}a_{45}a_{56}a_{63}\]
    记环路$L:M_3\to M_4\to M_5\to M_6\to M_3$，这四项分别对应以下四种情况：\begin{itemize}
        \item 不考虑节点$M_1,M_2$的自环，仅考虑环路L，符号为$-1$；
        \item 不考虑节点$M_2$的自环，考虑节点$M_1$的自环以及环路L，符号为$+1$；
        \item 不考虑节点$M_1$的自环，考虑节点$M_2$的自环以及环路L，符号为$+1$；
        \item 考虑节点$M_1,M_2$的自环以及环路L，符号为$-1$。
    \end{itemize}
\end{example}

最后，对于代数余子式$S^*_{n1}=\sum_{k}g_k\Delta_k$，我们考虑将$S$按第$n$行展开：
\[\det S=\Delta=\sum_{j=1}^{n}s_{nj}S^*_{nj}\]
这样，$S^*_{n1}$就对应$s_{n1}$的系数。行列式公式中某一项包含$s_{n1}$，说明置换$\sigma$中有$\sigma(1)=n$，即环路包含$M_1\to M_n$。将该环路的其余部分补全，就得到了一条从$M_1$到$M_n$的前向通路，其前向通路增益为$g_k$。去掉与该环路接触的节点后，剩下的部分正好对应$\Delta_k$。

注意，如果原图中不含有$M_1\to M_n$的支路，则$s_{n1}=0$，此时我们只好假想有一条“增益为0”的支路连接$M_1$和$M_n$，才能看出代数余子式的几何意义。因此，在梅森增益公式中，我们总是需要将前向通路考虑进来，而不管是否存在$M_1\to M_n$的支路。

\begin{example}
    考虑具有4个中间节点的信号流图：
    \begin{figure}[H]
        \centering
        \begin{tikzpicture}[auto, >=Stealth, node distance=2.2cm]
            % 节点（2mm空心小圆）
            \node[circle, draw, minimum size=2mm, inner sep=0pt, line width=0.6pt] (M1) {};
            \node[circle, draw, minimum size=2mm, inner sep=0pt, line width=0.6pt, right=of M1] (M2) {};
            \node[circle, draw, minimum size=2mm, inner sep=0pt, line width=0.6pt, right=of M2] (M3) {};
            \node[circle, draw, minimum size=2mm, inner sep=0pt, line width=0.6pt, right=of M3] (M4) {};
            % 节点标记，避免重叠
            \node[font=\small, below=3pt of M1] {$M_1$};
            \node[font=\small, below=3pt of M2] {$M_2$};
            \node[font=\small, below=3pt of M3] {$M_3$};
            \node[font=\small, below=3pt of M4] {$M_4$};
            % 边
            \draw[->] (M1) -- node[midway, above] {$a_{12}$} (M2);
            \draw[->] (M2) -- node[midway, above] {$a_{23}$} (M3);
            % a_{24}弯曲到M4
            \draw[->] (M2) to[out=-30,in=210,looseness=1.2] node[below right=1pt and -2pt, font=\small] {$a_{24}$} (M4);
            \draw[->] (M3) -- node[midway, above] {$a_{34}$} (M4);
            % a_{21}反馈支路弯曲
            \draw[->] (M2) to[out=150,in=30,looseness=1.3] node[above left=1pt and -2pt, font=\small] {$a_{21}$} (M1);
        \end{tikzpicture}
    \end{figure}
    仍设激励接入的节点为$M_1$，输出取自节点$M_n$，则前向通路有两条：
    \[M_1\to M_2\to M_3\to M_4,M_1\to M_2\to M_4\]
    我们来考察前向通路与系统矩阵的对应关系。增益矩阵为
    \[A=\begin{bmatrix}
        0 & a_{12} & 0 & 0 \\
        a_{21} & 0 & a_{23} & a_{24} \\
        0 & 0 & 0 & a_{34} \\
        0 & 0 & 0 & 0
    \end{bmatrix}\]
    因此系统矩阵为
    \[S=I-A=\begin{bmatrix}
        1 & s_{12} & 0 & 0 \\
        s_{21} & 1 & s_{23} & s_{24} \\
        0 & 0 & 1 & s_{34} \\
        0 & 0 & 0 & 1
    \end{bmatrix}\]
    
    这个系统并没有$M_4\to M_1$的支路，但我们假想有支路$a_{41}=-s_{41}=0$，这样就产生了环路
    $$M_1\to M_2\to M_3\to M_4\to M_1,M_1\to M_2\to M_4\to M_1$$
    前向通路增益分别为$g_1=a_{12}a_{23}a_{34}$和$g_2=a_{12}a_{24}$。

    计算特征行列式：原图中有环路$M_1\to M_2\to M_1$，环路增益为$L_1=a_{12}a_{21}$，特征行列式为$\Delta=1-L_1=1-a_{12}a_{21}$；对于两条前向通路，去掉所涉及的节点后不再有回路，因此$\Delta_1=\Delta_2=1$。根据梅森增益公式，有
    \[H=\frac{Y}{X}=\frac{g_1\Delta_1+g_2\Delta_2}{\Delta}=\frac{a_{12}a_{23}a_{34}+a_{12}a_{24}}{1-a_{12}a_{21}}\]
\end{example}

一般来说，我们直接用计算机分析信号流图表示的复杂系统，但在一些特定的信号流图中，使用几何意义更为方便。

\begin{example}[单入单出信号流图]
    系统的信号流图如图所示：
    \begin{figure}[H]
        \centering
        \begin{tikzpicture}[auto, >=Stealth, node distance=1.8cm, every node/.style={font=\small}]
            \tikzset{signode/.style={circle, draw, minimum size=2mm, inner sep=0pt, line width=0.6pt}, lab/.style={font=\small}}
            % nodes
            \node[signode] (X) {};
            \node[signode, right=2.2cm of X] (m1) {};
            \node[signode, below=1.2cm of m1,xshift=24mm] (m2) {};
            \node[signode, right=2.4cm of m1, yshift=12mm] (m3) {};
            \node[signode, right=2.2cm of m3] (Y) {}; 
            % labels
            \node[lab,left=1mm of X] {$X$};
            \node[lab,right=1mm of Y] {$Y$};
            \node[lab, above=0.5mm of m1] {$M_1$};
            \node[lab, below=0.5mm of m2] {$M_2$};
            \node[lab, above=0.5mm of m3] {$M_3$};
            % input connects only to m1 with unit gain
            \draw[->] (X) to node[midway, above, inner sep=1pt] {$1$} (m1);
            % output is taken from m3 with unit gain
            \draw[->] (m3) to node[midway, above, inner sep=1pt] {$1$} (Y);
            % complete directed subgraph among m1..m3 (H_ij for forward, G_ij for feedback)
            \draw[->] (m1) to[bend left=18] node[midway, above, inner sep=1pt] {$H_{12}$} (m2);
            \draw[->] (m2) to[bend left=18] node[midway, below, inner sep=1pt] {$G_{21}$} (m1);
            \draw[->] (m1) to[bend left=18] node[midway, above, inner sep=1pt,xshift=-6pt] {$H_{13}$} (m3);
            \draw[->] (m3) to[bend left=18] node[midway, below, inner sep=1pt,xshift=6pt] {$G_{31}$} (m1);
            \draw[->] (m2) to[bend left=18] node[midway, above, inner sep=1pt] {$H_{23}$} (m3);
            \draw[->] (m3) to[bend left=18] node[midway, below, inner sep=1pt] {$G_{32}$} (m2);
        \end{tikzpicture}
    \end{figure}
    \ding{172}使用矩阵进行分析：系统矩阵为
    \[\begin{bmatrix}
        1 & -G_{21} & -G_{31} \\
        -H_{12} & 1 & -G_{32} \\
        -H_{13} & -H_{23} & 1
    \end{bmatrix}\]
    特征行列式为
    $$\Delta=\begin{vmatrix}
        1 & -G_{21} & -G_{31} \\
        -H_{12} & 1 & -G_{32} \\
        -H_{13} & -H_{23} & 1
    \end{vmatrix}=1-G_{21}H_{12}-G_{31}H_{13}-G_{32}H_{23}-G_{21}G_{32}H_{13}-G_{31}H_{23}H_{12}$$
    在$\Delta\neq 0$时，系统矩阵可逆，其逆矩阵为\begin{align*}
        S^{-1}=\frac{1}{\Delta}\begin{bmatrix}
            1-G_{32}H_{23} & G_{21}+G_{31}H_{23} & G_{31}+G_{21}G_{32} \\
            H_{12}+H_{13}G_{32} & 1-G_{31}H_{13} & G_{32}+G_{31}H_{12} \\
            H_{13}+H_{12}H_{23} & H_{23}+G_{21}H_{13} & 1-G_{21}H_{12}
        \end{bmatrix}
    \end{align*}
    输出为$Y=M_3=\frac{1}{\Delta}(H_{13}+H_{12}H_{23})X$，因此系统增益为
    \[H=\frac{Y}{X}=\frac{H_{13}+H_{12}H_{23}}{\Delta}=\frac{H_{13}+H_{12}H_{23}}{1-G_{21}H_{12}-G_{31}H_{13}-G_{32}H_{23}-G_{21}G_{32}H_{13}-G_{31}H_{23}H_{12}}\]

    \ding{173}使用梅森增益公式的几何意义：图中共有五个环路，分别为\begin{align*}
        L_1&:M_1\to M_2\to M_1,\text{增益}L_1=G_{21}H_{12}\\
        L_2&:M_1\to M_3\to M_1,\text{增益}L_2=G_{31}H_{13}\\
        L_3&:M_2\to M_3\to M_2,\text{增益}L_3=G_{32}H_{23}\\
        L_4&:M_1\to M_2\to M_3\to M_1,\text{增益}L_4=G_{21}G_{32}H_{13}\\
        L_5&:M_1\to M_3\to M_2\to M_1,\text{增益}L_5=G_{31}H_{23}H_{12}
    \end{align*}
    这五个环路两两相交，因此特征行列式为
    \[\Delta=1-(L_1+L_2+L_3+L_4+L_5)=1-G_{21}H_{12}-G_{31}H_{13}-G_{32}H_{23}-G_{21}G_{32}H_{13}-G_{31}H_{23}H_{12}\]
    前向通路共有两条，增益分别为$g_1=H_{13}$和$g_2=H_{12}H_{23}$，去掉前向通路中的节点后没有环路，因此$\Delta_1=\Delta_2=1$。根据梅森增益公式，有
    \[H=\frac{g_1\Delta_1+g_2\Delta_2}{\Delta}=\frac{H_{13}+H_{12}H_{23}}{1-G_{21}H_{12}-G_{31}H_{13}-G_{32}H_{23}-G_{21}G_{32}H_{13}-G_{31}H_{23}H_{12}}\]
\end{example}

我们同样能够根据系统函数画出信号流图，但是一个的系统函数可以对应多种信号流图，例如两个系统的级联是可交换的。下面给出一种仅仅使用乘法器和积分器（系统函数为$s^{-1}$）将有理系统函数化为信号流图的标准方法，使用这种方法画出的信号流图是唯一的，并且具有相似的结构。

考虑如下信号流图：

\begin{center}
\scalebox{0.92}{%
\begin{tikzpicture}[>=Stealth, node distance=1.6cm, auto]
    % 节点：M1 为矩形，其余中间节点用直径为2的小圆圈表示，节点内部不显示编号
    \node (x) {$x(t)$};
    \node[draw, circle, minimum size=2mm, inner sep=0pt] (M1) [right=of x] {};
    \node[circle, draw, minimum size=2mm, inner sep=0pt] (M2) [right=of M1] {};
    \node[circle, draw, minimum size=2mm, inner sep=0pt] (M3) [right=of M2] {};
    \node (dots) [right=of M3] {$\cdots$};
    \node[circle, draw, minimum size=2mm, inner sep=0pt] (Mn_2) [right=of dots] {};
    \node[circle, draw, minimum size=2mm, inner sep=0pt] (Mn_1) [right=of Mn_2] {};
    \node[circle, draw, minimum size=2mm, inner sep=0pt] (Mn) [right=of Mn_1] {};
    \node (y) [right=of Mn] {$y(t)$};

    % 串联 s^{-1} 通路（链路上方标注，避免与其他箭头重叠）
    \draw[->] (x) -- (M1);
    \draw[->] (M1) -- node[above] {$s^{-1}$} (M2);
    \draw[->] (M2) -- node[above] {$s^{-1}$} (M3);
    \draw[->] (M3) -- node[above] {$s^{-1}$} (dots);
    \draw[->] (dots) -- node[below] {$s^{-1}$} (Mn_2);
    \draw[->] (Mn_2) -- node[below] {$s^{-1}$} (Mn_1);
    \draw[->] (Mn_1) -- node[below] {$b_n$} (Mn);
    \draw[->] (Mn) -- (y);

    % 直接到 Mn 的前向通路，标签放在上方并调整弯度避免重叠
    \draw[->, bend left=60] (M1) to node[above, yshift=1pt] {$b_0$} (Mn);
    \draw[->, bend left=50] (M2) to node[above, yshift=1pt] {$b_1$} (Mn);
    \draw[->, bend left=40] (M3) to node[above, yshift=1pt] {$b_2$} (Mn);
    \draw[->, bend left=30] (Mn_2) to node[above, yshift=1pt] {$b_{\,n-1}$} (Mn);

    % 反馈回到 M1（画在下方，标签置于弧线下方并微调，避免文字与箭头相交）
    \draw[->, bend left=30] (M2) to node[below, yshift=-1pt] {$-a_1$} (M1);
    \draw[->, bend left=40] (M3) to node[below, yshift=-1pt] {$-a_2$} (M1);
    \draw[->, bend left=50] (Mn_2) to node[below, yshift=-1pt] {$-a_{n-1}$} (M1);
    \draw[->, bend left=60] (Mn_1) to node[below, yshift=-1pt] {$-a_{n}$} (M1);
\end{tikzpicture}%
}
\end{center}

注意到这个信号流图的所有前向通路和环路均经过节点$M_1$，去掉节点$M_1$后不再有环路，因此不含$M_1$的子图的特征行列式$\Delta_k$均为1。利用梅森增益公式，可以计算出系统增益为
\begin{align*}
    H(s)&=\frac{Y(s)}{X(s)}=\frac{b_0 + b_1 s^{-1} + b_2 s^{-2} + \cdots + b_n s^{-n}}{1 + a_1 s^{-1} + a_2 s^{-2} + \cdots + a_n s^{-n}}
\end{align*}

可见，对于任意有理系统函数，都可以使用上述信号流图表示。注意系数$a_k,b_k$是可以为0的，因此即便分子分母不是相同度数的多项式，也可以使用该信号流图表示。

\subsection{最小相位系统}\label{sub:min_phase}

根据\ref{sub:stability_of_system}，对于因果且稳定的连续线性时不变系统，系统函数$H(s)$的所有极点均位于复平面的左半平面。如果我们还希望系统具有最小的相位延迟特性，从而具有最小的时延特性，则需要系统函数$H(s)$的所有零点也位于复平面的左半平面，这样的系统称为\textbf{最小相位系统}(minimum-phase system)。

\begin{definition}[最小相位系统]
    设因果且稳定的连续线性时不变系统的系统函数为$H(s)$，记$H(s)$的零点集为$\{z_k\}$，如果$\forall k,\text{Re}(z_k)<0$，则称该系统为\textbf{最小相位系统}。
\end{definition}

为了说明为什么这样的系统具有最小的延时特性，我们首先考虑一个例子：
\begin{example}
    设因果且稳定的连续线性时不变系统A，B的系统函数分别为
    \[H_A(s)=\frac{s+1}{s+2},H_B(s)=\frac{s-1}{s+2}\]
    这两个系统具有相同的幅频特性：
    \[|H_A(\mathrm{i}\omega)|=|H_B(\mathrm{i}\omega)|=\frac{\sqrt{\omega^2+1}}{\sqrt{\omega^2+4}}\]
    其中系统A为最小相位系统，系统B不是最小相位系统。计算两个系统的相频特性：
    \begin{align*}
        \varphi_A(\omega)&=\arg(\mathrm{i}\omega+1)-\arg(\mathrm{i}\omega+2)=\arctan(\omega)-\arctan\left(\frac{\omega}{2}\right)\\
        \varphi_B(\omega)&=\arg(\mathrm{i}\omega-1)-\arg(\mathrm{i}\omega+2)=-\arctan(\omega)-\arctan\left(\frac{\omega}{2}\right)
    \end{align*}

    在\ref{sub:distortion}中曾介绍，考察系统的延时特性，可以使用相时延$\tau_p$和群时延$\tau_g$。计算两个系统的相时延和群时延：
    \begin{align*}
        \tau_{pA}(\omega)&=-\frac{\varphi_A(\omega)}{\omega}=\frac{-\arctan(\omega)+\arctan\left(\frac{\omega}{2}\right)}{\omega}\\
        \tau_{pB}(\omega)&=-\frac{\varphi_B(\omega)}{\omega}=\frac{\arctan(\omega)+\arctan\left(\frac{\omega}{2}\right)}{\omega}\\
        \tau_{gA}(\omega)&=-\frac{d\varphi_A(\omega)}{d\omega}=-\frac{1}{1+\omega^2}+\frac{2}{4+\omega^2}\\
        \tau_{gB}(\omega)&=-\frac{d\varphi_B(\omega)}{d\omega}=\frac{1}{1+\omega^2}+\frac{2}{4+\omega^2}
    \end{align*}
    可以发现，
    \begin{align*}
        &\tau_{pA}(\omega)-\tau_{pB}(\omega)=-\frac{2\arctan(\omega)}{\omega}<0\\
        &\tau_{gA}(\omega)-\tau_{gB}(\omega)=-\frac{2}{1+\omega^2}<0
    \end{align*}
    也就是说，最小相位系统A的相时延和群时延均小于非最小相位系统B的相时延和群时延，因此系统A具有更好的延时特性。
\end{example}

\begin{proposition}
    具有相同幅频特性的因果且稳定的连续线性时不变系统中，最小相位系统具有最小的相时延和群时延，因此具有最佳的延时特性。
\end{proposition}
\begin{proof}
    我们对系统函数为有理函数的情况进行证明。设因果且稳定的连续线性时不变系统的系统函数为
    \[H(s)=K(s)\frac{\prod_{k=1}^{m}(s-z_k)}{\prod_{l=1}^{n}(s-p_l)}\]
    其中$K(s)$仅有左半平面内的零点而没有极点。$\text{Re}(p_l)<0,\forall l,\text{Re}(z_k)> 0,\forall k$。设系统函数$H(s)$的幅频特性为$|H(\mathrm{i}\omega)|=A(\omega)$，则
    \[A(\omega)=|K(s)|\frac{\prod_{k=1}^{m}\sqrt{\omega^2+|z_k|^2}}{\prod_{l=1}^{n}\sqrt{\omega^2+|p_l|^2}}\]
    设另一个系统的系统函数为
    \[\tilde{H}(s)=K(s)\frac{\prod_{k=1}^{m}(s-\tilde{z}_k)}{\prod_{l=1}^{n}(s-p_l)}\]
    其中$\tilde{z}_k=-\overline{z_k}$，即$z_k$关于虚轴的对称点，则$\tilde{H}(s)$的幅频特性与$H(s)$相同，且该系统是最小相位系统。计算两个系统的相频特性之差：
    \begin{align*}
        \varphi(\omega)-\tilde{\varphi}(\omega)&=\sum_{k=1}^{m}\left(\arg(\mathrm{i}\omega-z_k)-\arg(\mathrm{i}\omega-\tilde{z}_k)\right)\\
        &=\sum_{k=1}^{m}\left(\arctan\left(\frac{\omega-\text{Im}(z_k)}{-\text{Re}(z_k)}\right)-\arctan\left(\frac{\omega-\text{Im}(z_k)}{\text{Re}(z_k)}\right)\right)\\
        &=-2\sum_{k=1}^{m}\arctan\left(\frac{\omega-\text{Im}(z_k)}{\text{Re}(z_k)}\right)
    \end{align*}
    比较相时延和群时延：
    \begin{align*}
        \tau_p(\omega)-\tilde{\tau}_p(\omega)&=-\frac{\varphi(\omega)-\tilde{\varphi}(\omega)}{\omega}\\
        &=\frac{2}{\omega}\sum_{k=1}^{m}\arctan\left(\frac{\omega-\text{Im}(z_k)}{\text{Re}(z_k)}\right)>0\\
        \tau_g(\omega)-\tilde{\tau}_g(\omega)&=-\frac{d}{d\omega}(\varphi(\omega)-\tilde{\varphi}(\omega))\\
        &=2\sum_{k=1}^{m}\frac{\text{Re}(z_k)}{(\omega-\text{Im}(z_k))^2+\text{Re}(z_k)^2}>0
    \end{align*}
    因此，最小相位系统$\tilde{H}(s)$的相时延和群时延均小于非最小相位系统$H(s)$的相时延和群时延。
\end{proof}

需要指出，这个证明并不严谨。其一，这个证明没有说明为什么最小相位系统要求系统函数在虚轴上也不存在零点；其二，这个证明仅适用于有理函数。但是，更一般的情况需要用到哈代空间中函数的内外分解，我没能找到既严谨又通俗的证明，只能说这个证明能够提供一些直观的理解。

在最小相位系统的约束下，其\textbf{逆系统}也是因果且稳定的连续线性时不变系统。这是因为最小相位系统的系统函数$H(s)$的零点均位于复平面的左半平面，因此其倒数$H^{-1}(s)$的极点也均位于复平面的左半平面，从而保证了逆系统的因果性和稳定性。

\begin{example}
    设因果且稳定的连续线性时不变系统的系统函数及单位脉冲响应为
    \[H(s)=\frac{s+1}{s+2},h(t)=\delta(t)-\mathrm{e}^{-2t}u(t)\]
    则其逆系统的系统函数及单位脉冲响应为
    \[H^{-1}(s)=\frac{s+2}{s+1},h_{inv}(t)=\delta(t)+\mathrm{e}^{-t}u(t)\]
    该逆系统也是因果且稳定的连续线性时不变系统。将两个系统级联，其单位脉冲响应为
\begin{align*}
    h_{eq}(t)&=h(t)*h_{inv}(t)\\
    &=\delta(t)-\mathrm{e}^{-2t}u(t)+\mathrm{e}^{-t}u(t)-(\mathrm{e}^{-2t}u(t))*(\mathrm{e}^{-t}u(t))\\
    &=\delta(t)-\mathrm{e}^{-2t}u(t)+\mathrm{e}^{-t}u(t)-u(t)\int_0^t \mathrm{e}^{-2\tau}\mathrm{e}^{-(t-\tau)}d\tau\\
    &=\delta(t)
\end{align*}
\end{example}

对于逆系统，容易产生这样的误解：将两个系统级联，则系统函数为$H(s)\cdot H^{-1}(s)=1$，因此输出应与输入相同；考虑群时延和相时延，两个系统的群时延和相时延之和必定为0，因此只要两个系统中必有其中之一具有负的时延；进一步，考虑输入单频的正弦信号，在负的相时延下，输出信号应当提前出现，从而违反因果性。

事实上，两个系统都产生时延，与系统在级联后表现为恒等变换系统，二者并不矛盾，并且两个系统也的确满足因果性，上例中已经说明了这一点。

那么，以上“违反因果性”的结论在那一步出错了呢？实际上，这是因为单频正弦信号“提前出现”的表现，并没有破坏因果性，这是因为，当我们考虑系统的激励为正弦信号时，实际上假定这是一个\textbf{稳态}信号，即该正弦信号从$t=-\infty$时刻就已经存在了，我们观察到的相位“提前”，\textbf{实际上是相位延迟}，而级联所得的恒等变换系统总延迟为0，实际上是两个系统的相位延迟相互抵消了。因此，并不存在违反因果性的情况。

进一步，群时延之和为0，表现出调幅信号的包络峰提前出现的现象。这也没有违反因果性，因为响应起始时刻仍在激励之后，这种包络峰提前出现的现象，实际上是\textbf{相位失真引起的波形重构}。

\begin{example}[负的群时延]
    考虑调幅信号通过具有负的群时延的系统。设
    \[e(t)=\cos(\omega_m t)\cos(\omega_c t),H(s)=\frac{s+1}{s+2},h(t)=\delta-\mathrm{e}^{-2t}u(t)\]
    则系统的群时延为
    \[\tau_g(\omega)=-\frac{d\varphi(\omega)}{d\omega}=-\frac{1}{1+\omega^2}+\frac{2}{4+\omega^2}\]
    在低频段，群时延为负值。设基波频率$\omega_m=1$，载波频率$\omega_c=5$（这个载波的频率不算很高，这是为了使响应与激励区分开），则
    \begin{align*}
        e(t)=\cos(t)\cos(5t)=\frac{1}{2}[\cos(4t)+\cos(6t)],\tau_g(1)=-0.1
    \end{align*}
    由此可以预测，系统的响应中，包络峰将提前出现0.1s。

    计算系统的响应：
    \begin{align*}
        y(t)&=(h*e)(t)\\
        &=\cos(t)\cos(5t)-\frac{1}{2}u(t)\int_{0}^{t}\mathrm{e}^{-2\tau}[\cos(4(t-\tau))+\cos(6(t-\tau))]d\tau\\
        &=\cos(t)\cos(5t)-\frac{1}{2}u(t)\left[\frac{\cos(4t)+2\sin(4t)-\mathrm{e}^{-2t}}{10}+\frac{\cos(6t)+3\sin(6t)-\mathrm{e}^{-2t}}{20}\right]
    \end{align*}
    \begin{figure}[H]
        \centering
        \includegraphics[width=0.8\textwidth]{neg_taug}
        \caption{调幅信号通过具有负群时延系统的响应}
    \end{figure}
    上图中，依稀能辨认出-0.1s的群时延，但由于输入的是稳态信号，并没有违反因果性。
\end{example}

那么，对于暂态信号，情况又如何呢？考虑相时延
$$tau_p(\omega)=\dfrac{-\arctan(\omega)+\arctan(\frac{\omega}{2})}{\omega},\forall \omega>0,\tau_p(\omega)<0$$
即各频率分量均有负的相时延。

不过，时限函数不可能同时是带限的（见\ref{sub:time_freq_limit}），激励在全频域上均有频率分量，这当然可能使得响应的各分量提前出现，却在$t=0$时刻前抵消，最终呈现出响应与激励起始时刻相同，不违反因果性的结果。

\begin{example}[暂态信号经过逆系统]
    设
    $$e(t)=\chi_{[0,1]}(t),H(s)=\frac{s+1}{s+2}$$
    则系统的响应为
    \begin{align*}
        y(t)&=(h*\chi_{[0,1]})(t)\\
        &=\chi_{[0,1]}(t)-u(t)\int_0^t \chi_{[0,1]}(\tau)\mathrm{e}^{-2(t-\tau)}d\tau
    \end{align*}
    单位脉冲响应$h(t)$是因果的，因此响应也是因果的。

    在$t<1$时，相应为\begin{align*}
        y(t)&=\chi_{[0,1]}(t)-u(t)\int_0^t \mathrm{e}^{-2(t-\tau)}d\tau\\
        &=\chi_{[0,1]}(t)-\frac{1}{2}u(t)(1-\mathrm{e}^{-2t})
    \end{align*}
    在$t\geq 1$时，相应为\begin{align*}
        y(t)&=\chi_{[0,1]}(t)-u(t)\int_0^1 \mathrm{e}^{-2(t-\tau)}d\tau\\
        &=-\frac{1}{2}u(t)(\mathrm{e}^{-2t}-\mathrm{e}^{-2(t-1)})
    \end{align*}
    \begin{figure}[H]
        \centering
        \includegraphics[width=0.8\textwidth]{inv_system}
        \caption{暂态信号通过逆系统的响应}
    \end{figure}
    响应与激励的起始时刻一致，满足因果性。
\end{example}
